% !TeX root = ../../../main.tex
% !TeX encoding = UTF-8

% ===== BEGIN ORIGINAL TYPESETTING TESTS: delete this whole block later =====
\subsection*{原创排版测试（非教材原题）}
以下内容仅用于测试环境、编号及引用，不计入教材习题完成进度。
其题号中的 T 表示测试；手动题号与原书无关。

\begin{definition}[距离的测试记号]
\label{def:b1-c1-s1-test-distance}
对 $x,y\in\R$，记 $d(x,y)=\abs{x-y}$。
\end{definition}

\begin{lemma}[三角不等式测试]
\label{lem:b1-c1-s1-test-triangle}
对任意 $a,b\in\R$，有 $\abs{a+b}\leq\abs{a}+\abs{b}$。
\end{lemma}
\begin{proof}
由 $-\abs{a}\leq a\leq\abs{a}$ 和
$-\abs{b}\leq b\leq\abs{b}$ 相加即得结论。
\end{proof}

\begin{exercise}[T1]
\label{ex:b1-c1-s1-test-reverse}
设 $x,y\in\R$，证明
\begin{equation}
  \label{eq:b1-c1-s1-test-reverse}
  \abs{\abs{x}-\abs{y}}\leq\abs{x-y}.
\end{equation}
\end{exercise}
\begin{solution}
由引理~\ref{lem:b1-c1-s1-test-triangle}，
$\abs{x}=\abs{(x-y)+y}\leq\abs{x-y}+\abs{y}$，
故 $\abs{x}-\abs{y}\leq\abs{x-y}$。
交换 $x,y$ 可得反向估计，从而得到式~\eqref{eq:b1-c1-s1-test-reverse}。
\end{solution}
\begin{idea}
分别估计差与其相反数，便可处理外层绝对值。
\end{idea}
\begin{alternative}
不妨设 $\abs{x}\geq\abs{y}$，直接对 $x=(x-y)+y$ 使用三角不等式；
另一种情况由对称性得到。
\end{alternative}
\begin{generalization}
任意赋范空间中的 $u,v$ 都满足
$\abs{\norm{u}-\norm{v}}\leq\norm{u-v}$，证明完全相同。
\end{generalization}

\begin{exercise}
\label{ex:b1-c1-s1-test-limit}
（自动编号测试）设实数列 $x_n\to a$，证明 $\abs{x_n}\to\abs{a}$。
\end{exercise}
\begin{solution}
任取 $\eps>0$。存在正整数 $N$，使得对所有 $n\geq N$，
有 $\abs{x_n-a}<\eps$。应用习题~\ref{ex:b1-c1-s1-test-reverse}，得到
\[
  \abs{\abs{x_n}-\abs{a}}\leq\abs{x_n-a}<\eps.
\]
这正是所需收敛结论的定义。
\end{solution}

\begin{theorem}[连续性的测试结论]
\label{thm:b1-c1-s1-test-continuity}
函数 $f\colon\R\to\R$，$f(x)=\abs{x}$ 在每一点连续。
\end{theorem}
\begin{proof}
给定 $a\in\R$ 及 $\eps>0$，取 $\delta=\eps$，再应用
式~\eqref{eq:b1-c1-s1-test-reverse} 即得。
\end{proof}

\begin{proposition}[复合的测试结论]
若实函数 $g$ 在 $a$ 处连续，则 $\abs{g}$ 在 $a$ 处连续。
\end{proposition}
\begin{proof}
由 $\abs{\abs{g(x)}-\abs{g(a)}}\leq\abs{g(x)-g(a)}$ 及连续性定义即得。
\end{proof}

\begin{corollary}
若 $x_n\to 0$，则 $\abs{x_n}\to 0$。
\end{corollary}
\begin{proof}
在习题~\ref{ex:b1-c1-s1-test-limit} 中令 $a=0$ 即得。
\end{proof}

\begin{example}[命令测试]
有限集合 $A=\set{0,1}\subset\N\subset\Z\subset\Q\subset\R\subset\C$，
且 $\int_0^1 x\dd x=\tfrac12$。
此处采用 $0\in\N$ 的约定；正式题解应按所用教材明确约定。
\end{example}

\begin{remark}[交叉引用测试]
本节属于第~\ref{chap:b1-c1}~章的第~\ref{sec:b1-c1-s1}~节。
定义~\ref{def:b1-c1-s1-test-distance} 与
定理~\ref{thm:b1-c1-s1-test-continuity} 的引用可点击跳转。
其他章节也使用相同的全局标签机制。
\end{remark}
% ===== END ORIGINAL TYPESETTING TESTS =====
